% GENERATED FILE. Edit canonical lessons, then run npm run generate:books.
% canonical-source-hash: fnv1a64:0c8dc57227d9907c
% canonical-lessons: ES-C307-esposa, ES-C307-tio, ES-C307-tia, ES-C307-primo, ES-C307-prima

\chapter{The Wider Circle --- Wife, Uncle, Aunt, Cousin}
\label{ch:esposos-y-tios}
\hlchaptermodality{307}

\begin{chapteropening}
\textbf{By the end of this chapter:} \emph{I can name a wife, an uncle, an aunt and a cousin in Spanish, in either gender, and I know why Spanish and English chose different ancient words for the same relatives.}

Complete ES-C307-prima: name the wider circle of relatives and pair each -o form with its -a form.
\end{chapteropening}

\section[la esposa]{la esposa --- the pledge that binds}

\label{lesson:ES-C307-esposa}



\begin{quote}
Two from last time, before anything new is added.
\end{quote}

\begin{sounds}
\begin{itemize}
\raggedright
  \item Say it once slowly, then at speed. $\to$ pronunciation reference
\end{itemize}
\end{sounds}

\begin{cousinweb}
\textbf{la esposa} --- \textquotedblleft{}the wife.\textquotedblright{}

The same pledge, the same ending switch. English \emph{spouse} does not bother to mark who it means; Spanish always does.

Now the strange part, and it is worth knowing because it is unmistakable once you have heard it. In the plural, \textbf{las esposas} also means \emph{handcuffs}. The metaphor is not a joke about marriage --- it runs straight out of the root. A pledge is a thing that binds you, so the binding thing takes the pledge's name.

English does the identical move and nobody blinks: you are \emph{bound} by a promise, you enter into \emph{bonds}, and a \emph{wedlock} is a lock. Spanish has aimed the word at the object, English at the feeling.
\end{cousinweb}

\begin{grammarlens}[title={What you've built}]
Six words in, and the -o/-a switch has done the work of six separate nouns.
\end{grammarlens}

\subsection*{Your turn}
Say these aloud:

\begin{itemize}
\raggedright
  \item \emph{la esposa}
  \item \emph{el esposo y la esposa}
  \item \emph{la esposa y la hija}
\end{itemize}

\begin{tcolorbox}[breakable,colback=teal!4,colframe=teal!35!black,title={Before you move on}]
What does \emph{las esposas} mean besides wives? (Handcuffs.) Why is that not a joke? (A pledge binds --- the root means to pledge.) Which English word does the same thing? (\emph{Bound}, or \emph{wedlock}.)
\end{tcolorbox}
\section[el tío]{el tío --- the Greek one}

\label{lesson:ES-C307-tio}



\begin{quote}
Two from last time, before anything new is added.
\end{quote}

\begin{sounds}
\begin{itemize}
\raggedright
  \item Say it once slowly, then at speed. $\to$ pronunciation reference
\end{itemize}
\end{sounds}

\begin{cousinweb}
\textbf{el tío} --- \textquotedblleft{}the uncle.\textquotedblright{}

This one is not Latin at its source. Greek \textbf{theîos} meant uncle, Late Latin borrowed it as \textbf{thius}, and Spanish carried it down to \emph{tío}.

Here is why that matters to you. Two lessons ago you learned that English \emph{uncle} comes from Latin \emph{avunculus}, the \textquotedblleft{}little grandfather\textquotedblright{} inside \emph{abuelo}. So the two languages had both words available and each kept the other one. English went Latin; Spanish went Greek. Same relative, two different ancestries, and neither word looks anything like the other.

The written accent on the \textbf{í} is doing real work. It splits \emph{tí-o} into two beats and keeps the stress on the first. Without it the word would collapse into one syllable.
\end{cousinweb}

\begin{grammarlens}[title={What you've built}]
A relative English and Spanish named from opposite ends of the ancient world.
\end{grammarlens}

\subsection*{Your turn}
Say these aloud:

\begin{itemize}
\raggedright
  \item \emph{el tío}
  \item \emph{el tío y el abuelo}
  \item \emph{el tío, señor}
\end{itemize}

\begin{tcolorbox}[breakable,colback=teal!4,colframe=teal!35!black,title={Before you move on}]
Where does \emph{tío} come from? (Greek \emph{theîos}.) Where does English \emph{uncle} come from? (Latin \emph{avunculus} --- the \emph{abuelo} root.) What is the accent on the í doing? (Splitting two beats, holding the stress.)
\end{tcolorbox}
\section[la tía]{la tía --- a third word for the same relative}

\label{lesson:ES-C307-tia}



\begin{quote}
Two from last time, before anything new is added.
\end{quote}

\begin{sounds}
\begin{itemize}
\raggedright
  \item Say it once slowly, then at speed. $\to$ pronunciation reference
\end{itemize}
\end{sounds}

\begin{cousinweb}
\textbf{la tía} --- \textquotedblleft{}the aunt.\textquotedblright{}

Spanish swings the Greek word feminine and is done. English could not: \emph{aunt} comes from Latin \textbf{amita}, the father's sister, which is neither \emph{thius} nor \emph{avunculus}. That is three separate ancient words in play for two relatives, and the two languages picked differently every time.

The accent works exactly as it did in \emph{tío} --- \emph{tí-a}, two beats, stress on the first.

One thing to be ready for. In everyday speech across much of the Spanish-speaking world, \emph{tío} and \emph{tía} get used for people who are not relatives at all, the way English uses \emph{dude} or \emph{mate}. That is casual, not neutral; you will hear it long before you need it.
\end{cousinweb}

\begin{grammarlens}[title={What you've built}]
Eight words, and a clear picture of two languages choosing different ancestors.
\end{grammarlens}

\subsection*{Your turn}
Say these aloud:

\begin{itemize}
\raggedright
  \item \emph{la tía}
  \item \emph{el tío y la tía}
  \item \emph{la tía y la abuela}
\end{itemize}

\begin{tcolorbox}[breakable,colback=teal!4,colframe=teal!35!black,title={Before you move on}]
Where does English \emph{aunt} come from? (Latin \emph{amita}.) How many ancient words are in play for uncle and aunt? (Three --- \emph{thius}, \emph{avunculus}, \emph{amita}.) Is calling a stranger \emph{tía} neutral? (No --- casual.)
\end{tcolorbox}
\section[el primo]{el primo --- the one that just means first}

\label{lesson:ES-C307-primo}



\begin{quote}
Two from last time, before anything new is added.
\end{quote}

\begin{sounds}
\begin{itemize}
\raggedright
  \item Say it once slowly, then at speed. $\to$ pronunciation reference
\end{itemize}
\end{sounds}

\begin{cousinweb}
\textbf{el primo} --- \textquotedblleft{}the cousin.\textquotedblright{}

You have met \textbf{primus} before, in \emph{primero}. It means \textquotedblleft{}first,\textquotedblright{} and it is the root under English \emph{prime}, \emph{primary}, \emph{premier} and \emph{primitive}.

So why is a cousin called \textquotedblleft{}first\textquotedblright{}? Because Latin said \textbf{consobrinus primus} --- \textquotedblleft{}first cousin,\textquotedblright{} exactly as English still does. Spanish dropped the long word and kept the adjective. What is left is a relative whose entire name is a ranking: first.

English did the reverse and kept the long word. \emph{Cousin} is worn down from \emph{consobrinus}, which is \emph{con-} plus \emph{sobrinus}, \textquotedblleft{}sister's child.\textquotedblright{} Both languages took half of the same Latin phrase --- and they took different halves.
\end{cousinweb}

\begin{grammarlens}[title={What you've built}]
An old root paying off again: \emph{primero} and \emph{primo} are the same word.
\end{grammarlens}

\subsection*{Your turn}
Say these aloud:

\begin{itemize}
\raggedright
  \item \emph{el primo}
  \item \emph{el primo y el tío}
  \item \emph{el nombre del primo}
\end{itemize}

\begin{tcolorbox}[breakable,colback=teal!4,colframe=teal!35!black,title={Before you move on}]
What does \emph{primus} mean? (First.) Which word did you already have from it? (\emph{Primero}.) Which half of \emph{consobrinus primus} did English keep? (\emph{Consobrinus}, which became \emph{cousin}.)
\end{tcolorbox}
\section[la prima]{la prima --- first lady, first face}

\label{lesson:ES-C307-prima}



\begin{quote}
Two from last time, before anything new is added.
\end{quote}

\begin{sounds}
\begin{itemize}
\raggedright
  \item Say it once slowly, then at speed. $\to$ pronunciation reference
\end{itemize}
\end{sounds}

\begin{cousinweb}
\textbf{la prima} --- \textquotedblleft{}the cousin,\textquotedblright{} of a girl or woman.

The ending switch again. What makes this one easy is that English already says \emph{prima} out loud without translating it.

A \textbf{prima donna} is the \textquotedblleft{}first lady\textquotedblright{} of an opera company --- the leading singer. The phrase turned into an insult later, but the words are plain: first woman.

\textbf{Prima facie} evidence is evidence that holds up \textquotedblleft{}at first face,\textquotedblright{} on first look. Same \emph{prima}, same meaning, borrowed whole out of Latin and never naturalised.

So you already knew this word. You knew it as a label for opera singers and legal arguments, and it is the same word you now use for your cousin.
\end{cousinweb}

\begin{grammarlens}[title={What you've built}]
Ten words --- and the chapter's last. Two families of relatives, named.
\end{grammarlens}

\subsection*{Your turn}
Say these aloud:

\begin{itemize}
\raggedright
  \item \emph{la prima}
  \item \emph{el primo y la prima}
  \item \emph{la prima y la tía}
\end{itemize}

\begin{tcolorbox}[breakable,colback=teal!4,colframe=teal!35!black,title={Before you move on}]
What is a \emph{prima donna} literally? (First lady.) What is \emph{prima facie}? (At first face.) What is a \emph{prima} to you now? (A cousin.)
\end{tcolorbox}
